\section{Related Work}

\paragraph{What does a learned solver contribute?}
DIMES meta-learns the initialization of a distribution over discrete solutions, then samples feasible sets from its vertex scores \citep{qiu2022dimes}. Graph convolutional networks can predict vertex scores for tree search \citep{li2018guided}, and a policy can defer difficult vertex assignments during construction \citep{ahn2020defer}. Both decisions interact with the search that consumes them. In a reimplementation of guided MIS tree search, \citet{boether2022tree} found that kernelization and local search could preserve solution quality even with random scores. \citet{wu2026unrealized} compared neural MIS solvers with single-core KaMIS and degree-based greedy construction, and found the classical methods competitive even on the neural solvers' training distributions. ML4CO-Bench-101 evaluates the contributions of learned components and surrounding heuristics across graph optimization problems \citep{ma2025ml4co}. Our comparison fixes the greedy starting set, exchanges and branching rules, then charges each ranking for its own inference time. This measures what a learned priority adds to the set a solver can actually deliver.

\paragraph{How does a candidate become an answer?}
OnlineMIS applies graph reductions during local search \citep{dahlum2016accelerating}. Its search can retain a feasible incumbent while exploring. Neural samplers can instead move through conflicting states. Regularized Langevin dynamics controls the expected distance of a transition \citep{feng2025regularized}, and graph diffusion solvers convert denoised states into discrete solutions \citep{sun2023difusco,sanokowski2024diffusion}. For these methods both the search trajectory and the cost of recovering a feasible set matter. The primal integral captures how an incumbent's objective gap evolves through time \citep{berthold2013primal}. We measure the feasible cardinality at a specified interruption time, when its vertices first become available on the host, and when the solve returns. These observations can be made even when the optimum is unknown.
